\documentclass[../../main.tex]{subfiles}
\graphicspath{{\subfix{../../image/}}}

\begin{document}

\section{习题}\label{section:ch04-exercises}

在以下习题中，除非另作说明，所有给定函数均设为光滑。

\begin{exercise}\label{exercise:ch04-01}
用分离变量法求 PDE
\begin{align*}
u_{x_1}^2u_{x_1x_1}+2u_{x_1}u_{x_2}u_{x_1x_2}+u_{x_2}^2u_{x_2x_2}=0
\end{align*}
在 $\mathbb{R}^2$ 中的一个非平凡解。
\end{exercise}

\begin{exercise}\label{exercise:ch04-02}
考虑 $\mathbb{R}^2$ 中 Laplace 方程 $\Delta u=0$，带 Cauchy 数据 $u=0$、$\dfrac{\partial u}{\partial x_2}=\dfrac{1}{n}\sin(nx_1)$ 在 $\{x_2=0\}$ 上。用分离变量法导出解 $u=\dfrac{1}{n^2}\sin(nx_1)\sinh(nx_2)$。当 $n\to\infty$ 时 $u$ 如何？Laplace 方程的 Cauchy 问题是否适定？（本例由 Hadamard 给出。）
\end{exercise}

\begin{exercise}\label{exercise:ch04-03}
求 $v$ 与 $\sigma$ 的显式公式，使 $u(x,t)\triangleq v(x-\sigma t)$ 是非线性扩散方程 $u_t-u_{xx}=f(u)$ 的行波解，其中 $f(z)=-2z^3+3z^2-z$。设 $\lim_{s\to\infty}v=1$、$\lim_{s\to-\infty}v=0$、$\lim_{s\to\pm\infty}v'=0$。（提示：用 $v'$ 乘方程 $v''+\sigma v'+f(v)=0$ 并积分，以确定 $\sigma$ 的值。）
\end{exercise}

\begin{exercise}\label{exercise:ch04-04}
若求非线性椭圆方程 $-\Delta u=u^p$ 在 $\mathbb{R}^n$ 中的径向解 $u(x)=v(r)$（$r=|x|$、$p>1$），则导出非自治 ODE $v''+\dfrac{n-1}{r}v'+v^p=0$。证明 Emden--Fowler 变换 $t\triangleq\log r$、$x(t)\triangleq e^{\frac{2t}{p-1}}v(e^t)$ 把它化为新未知量 $x=x(t)$ 的一个自治 ODE。
\end{exercise}

\begin{exercise}\label{exercise:ch04-05}
求非线性热方程 $u_t-\Delta(u^\gamma)=0$ 的一个非负伸缩不变解，形如 $u(x,t)=t^{-\alpha}v(xt^{-\beta})$，其中 $\dfrac{n-2}{n}<\gamma<1$。解应当 $|x|\to\infty$ 时代数趋于零。
\end{exercise}

\begin{exercise}\label{exercise:ch04-06}
求 $-\Delta u+u^{\frac{n+2}{n-2}}=0$ 在 $B(0,1)$ 中、形如 $u=\alpha(1-|x|^2)^{-\beta}$（$\alpha,\beta>0$）的解。本例表明非线性 PDE 的解可在区域内有限而在其边界上处处趋于无穷。
\end{exercise}

\begin{exercise}\label{exercise:ch04-07}
考虑黏性守恒律 $u_t+F(u)_x-au_{xx}=0$ 在 $\mathbb{R}\times(0,\infty)$ 中，$a>0$、$F$ 一致凸。
\begin{enumerate}[(a)]
\item 证明：若 $u(x,t)=v(x-\sigma t)$ 且 $v$ 由公式 $s=\displaystyle\int_{v(s)}^c\frac{a}{F(z)-\sigma z+b}\,\mathrm{d}z$ 隐式定义（$b,c$ 为常数），则 $u$ 解此方程。
\item 证明：当且仅当 $\sigma=\dfrac{F(u_l)-F(u_r)}{u_l-u_r}$（$u_l>u_r$）时，存在满足 $\lim_{s\to-\infty}v(s)=u_l$、$\lim_{s\to\infty}v(s)=u_r$ 的行波。
\item 设 $u^\varepsilon$ 为 $a=\varepsilon$ 时上述行波解、$u^\varepsilon(0,0)=\dfrac{u_l+u_r}{2}$。计算 $\lim_{\varepsilon\to 0}u^\varepsilon$ 并解释。
\end{enumerate}
\end{exercise}

\begin{exercise}\label{exercise:ch04-08}
证明：若 $u$ 是 \S4.3 Schrödinger 方程问题由公式\eqref{align:ch04-schr-formula} 给出的解，则 $\|u(\cdot,t)\|_{L^\infty(\mathbb{R}^n)}\leqslant\dfrac{1}{(4\pi|t|)^{n/2}}\|g\|_{L^1(\mathbb{R}^n)}$ 对每个 $t\neq 0$。
\end{exercise}

\begin{exercise}\label{exercise:ch04-09}
设 $u$ 解非线性热方程 $u_t=\dfrac{u_{xx}}{u_x^2}$ 在 $\mathbb{R}\times(0,\infty)$ 中、$u_x>0$。令 $v$ 为每个时刻 $t>0$ 时 $u$ 在变量 $x$ 中的反函数，即 $y=u(x,t)\iff x=v(y,t)$。证明 $v$ 解一个线性 PDE。
\end{exercise}

\begin{exercise}\label{exercise:ch04-10}
求函数 $f:\mathbb{R}^3\to\mathbb{R}$、$f=f(z,p_1,p_2)$，使若 $u$ 是旋转波动方程 $u_{xt}=0$ 的任一解，则 $w\triangleq f(u,u_x,u_t)$ 解 Liouville 方程 $w_{xt}=e^w$。（提示：证明 $f$ 必具形式 $f(z,p_1,p_2)=a(z)+b(p_1)+c(p_2)$。）
\end{exercise}

\begin{exercise}[Lax 对]\label{exercise:ch04-11}设 $\{L(t)\}_{t\geqslant 0}$ 是某实 Hilbert 空间 $H$ 上的对称线性算子族，满足演化方程 $\dot L=[B,L]=BL-LB$。又设有相应特征值族 $\{\lambda(t)\}$ 与特征向量族 $\{w(t)\}$：$L(t)w(t)=\lambda(t)w(t)$。设 $L,B,\lambda,w$ 均光滑依赖于 $t$。证明 $\dot\lambda\equiv 0$。（提示：对 $t$ 微分恒等式 $Lw=\lambda w$，计算 $(\dot\lambda w,w)$。）
\end{exercise}

\begin{exercise}[续上]\label{exercise:ch04-12}给定 $u=u(x,t)$，定义线性算子 $L(t)v\triangleq-v_{xx}+uv$、$B(t)v=-4v_{xxx}+6uv_x+3u_xv$。证明 $(\dot L-[B,L])v=(u_t+u_{xxx}-6uu_x)v$。因而若 $u$ 解 KdV 方程 $u_t-6uu_x+u_{xxx}=0$，则线性算子 $L(t)$ 的特征值不随时间改变。
\end{exercise}

\begin{exercise}\label{exercise:ch04-13}
证明可从\cref{example:ch04-singular-transport}的初值问题构造显式解，形如 $v(z,t)=\dfrac{1}{\sigma(t)}\dfrac{1}{(\pi\gamma(t))^{1/2}}e^{-z^2/\gamma(t)}$（$z\in\mathbb{R}$、$t>0$），函数 $\gamma(t)$ 待定。代入 PDE 确定 $\gamma$ 应满足的 ODE。此 ODE 的初值是什么？
\end{exercise}

\begin{exercise}\label{exercise:ch04-14}
设 $u^\varepsilon,v^\varepsilon$ 解组
\begin{align*}
\begin{cases}
u_t^\varepsilon+\dfrac{1}{\varepsilon}u_x^\varepsilon=\dfrac{(v^\varepsilon)^2-(u^\varepsilon)^2}{\varepsilon^2},\\
v_t^\varepsilon-\dfrac{1}{\varepsilon}v_x^\varepsilon=\dfrac{(u^\varepsilon)^2-(v^\varepsilon)^2}{\varepsilon^2}.
\end{cases}
\end{align*}
设可写 $u^\varepsilon=u_0+\varepsilon u_1+\varepsilon u_2^\varepsilon$、$v^\varepsilon=v_0+\varepsilon v_1+\varepsilon v_2^\varepsilon$，其中 $u_0,v_0,u_1,v_1$ 光滑、$u_0,v_0>0$，且 $u_2^\varepsilon,v_2^\varepsilon$ 及其导数关于 $\varepsilon$ 一致有界。证明 $u_0\equiv v_0\triangleq w$，且 $w$ 解非线性热方程 $w_t-\dfrac{1}{4}(\log w)_{xx}=0$。
\end{exercise}

\begin{exercise}\label{exercise:ch04-15}
在 \S4.5.3 引理\cref{lemma:ch04-stationary-quadratic} 的证明中，论证把 $e^{-z^2}$ 沿直线 $\Gamma$ 的积分变形为沿实轴的积分。
\end{exercise}

\begin{exercise}\label{exercise:ch04-16}
用 \S4.5.3 引理\cref{lemma:ch04-stationary-quadratic} 讨论：由公式\eqref{align:ch04-schr-formula} 定义的 $u$ 在何种意义下当 $t\to 0^+$ 收敛到初值 $g$。
\end{exercise}

\begin{exercise}\label{exercise:ch04-17}
设 $n=1$，$u^\varepsilon$ 解 $-(a(x/\varepsilon)u_x^\varepsilon)_x=f$ 在 $(0,1)$ 中、$u^\varepsilon(0)=u^\varepsilon(1)=0$，其中 $a$ 光滑正、1-周期，$f\in L^2(0,1)$。
\begin{enumerate}[(a)]
\item 证明 $u^\varepsilon\rightharpoonup u$ 弱于 $H_0^1(0,1)$，其中 $u$ 解 $-\overline{a}u_{xx}=f$、$u(0)=u(1)=0$，$\overline{a}\triangleq\left(\int_0^1 a(y)^{-1}\,\mathrm{d}y\right)^{-1}$。
\item 验证此答案与 \S4.5.4 的均匀化结论一致。
\end{enumerate}
\end{exercise}

\begin{exercise}[均匀化中的变分原理]\label{exercise:ch04-18}设 $A(y)=((a^{ij}(y)))$ 对称正定、$Q$-周期。
\begin{enumerate}[(a)]
\item 导出对每个 $\xi\in\mathbb{R}^n$ 的变分公式 $\xi\cdot\overline{A}\xi=\min_w\left\{\int_Q Dw\cdot A(y)Dw\,\mathrm{d}y\mid w=y\cdot\xi+v,\ v\ Q\text{-周期}\right\}$。
\item 导出对偶变分公式 $\eta\cdot\overline{A}^{-1}\eta=\min_{\boldsymbol{\sigma}}\left\{\int_Q\boldsymbol{\sigma}\cdot A(y)^{-1}\boldsymbol{\sigma}\,\mathrm{d}y\mid\int_Q\boldsymbol{\sigma}\,\mathrm{d}y=\eta,\ \operatorname{div}\boldsymbol{\sigma}=0,\ \boldsymbol{\sigma}\ Q\text{-周期}\right\}$。
\item 证明 $\left(\int_Q A(y)^{-1}\,\mathrm{d}y\right)^{-1}\leqslant\overline{A}\leqslant\int_Q A(y)\,\mathrm{d}y$。
\end{enumerate}
\end{exercise}

\begin{exercise}\label{exercise:ch04-19}
证明直线 $\{t=0\}$ 对热方程 $u_t=u_{xx}$ 是特征的。证明不存在热方程在 $\mathbb{R}\times\mathbb{R}$ 中的解析解 $u$，使 $u=\dfrac{1}{1+x^2}$ 在 $\{t=0\}$ 上。（提示：设有解析解，算其系数，证明所得幂级数除 $(0,0)$ 外发散。本例由 Kovalevskaya 给出。）
\end{exercise}

\section{注释与参考文献}\label{section:ch04-references}

\begin{itemize}
\item[\S4.1] 分离变量法更多内容见 Pinsky \cite{P}、Strauss \cite{St2}、Thoe--Zachmanoglou \cite{T-Z}、Weinberger \cite{We}。Turing 不稳定见 Grindrod \cite{Gr}、Nishiura \cite{N}。
\item[\S4.2] 非线性色散 PDE 理论见 Linares--Ponce \cite{L-P}、Tao \cite{To}。双稳态方程行波的讨论由 C. Jones 提供，Barenblatt 解的推导由 J.-L.\ Vazquez 告知。P. Olver 的书 \cite{O1} 详述 PDE 的对称方法。
\item[\S4.3] Fourier 变换技术更多内容见 Stein--Weiss \cite{S-W}、Stein \cite{Se}、Hörmander \cite{H}、Rauch \cite{R1}、Treves \cite{T}、Zworski \cite{Zi}。定理\cref{theorem:1222-Fourier变换的性质} 的证明由 Y. Yu 改进，Schrödinger 方程由 M. Weinstein 协助。能量渐近均分公式\eqref{align:ch04-energy-split} 基于 Brodsky。例\cref{example:ch04-telegraph} 来自 Pinsky \cite{P}。Radon 变换适用范围的更精细研究见 Helgason \cite{Hg}，亦见 John \cite{J1}、Lax \cite{Lx3}。
\item[\S4.4] hodograph 与 Legendre 变换见 Courant--Hilbert \cite{C-H}。
\item[\S4.5] \S4.5.1 由 J. Neu 贡献。\S4.5.3 基于 J. Ralston 的课堂讲义，追随 Hörmander。\S4.5.4 均匀化的讨论追随 Bensoussan--Lions--Papanicolaou \cite{B-L-P}。摄动理论见 Neu \cite{Ne2}，几何光学见 Rauch \cite{R2}。
\item[\S4.6] 见 Folland \cite{F1} 第 1 章、John \cite{J2} 第 3 章、DiBenedetto \cite{DB1} 第 1 章。
\item[\S4.7] 习题\cref{exercise:ch04-09}、\cref{exercise:ch04-10} 基于 Zwillinger \cite{Zw} 中的例子。习题\cref{exercise:ch04-18} 由 R. Kohn 贡献。习题\cref{exercise:ch04-19} 来自 Mikhailov \cite{M}。
\end{itemize}

\end{document}
